\UnnumberedChapter{LỜI MỞ ĐẦU}

Trong bối cảnh chuyển đổi số đang diễn ra mạnh mẽ trong lĩnh vực giáo dục, việc xây dựng các hệ thống quản lý thông minh nhằm nâng cao hiệu quả vận hành và hỗ trợ ra quyết định đang trở thành xu hướng tất yếu. Mô hình **Smart Campus** được xem là một giải pháp hiện đại, cho phép kết nối và quản lý tập trung nhiều dịch vụ trong khuôn viên trường học thông qua các công nghệ như Internet of Things (IoT), điện toán đám mây, kiến trúc Microservices và các giao diện lập trình ứng dụng (API).

Trong một hệ thống Smart Campus, việc phát hiện và xử lý kịp thời các sự kiện bất thường có vai trò đặc biệt quan trọng. Những cảnh báo liên quan đến an ninh, môi trường hoặc thiết bị cần được truyền tải nhanh chóng đến người quản trị nhằm giảm thiểu rủi ro và nâng cao khả năng phản ứng của hệ thống. Chính vì vậy, việc xây dựng một dịch vụ chuyên trách tiếp nhận và gửi thông báo là yêu cầu cần thiết trong kiến trúc tổng thể của hệ thống.

Xuất phát từ yêu cầu đó, nhóm B7 thực hiện đề tài **"Xây dựng Notification Service cho hệ thống Smart Campus Operations Platform"** trong khuôn khổ học phần **Dịch vụ kết nối và Công nghệ nền tảng (FIT4110)**. Notification Service được xây dựng theo kiến trúc Microservices với nhiệm vụ tiếp nhận các cảnh báo từ Core Business Service, gửi thông báo đến người dùng thông qua Telegram Bot, lưu trữ lịch sử xử lý vào cơ sở dữ liệu PostgreSQL và cung cấp dữ liệu cho Analytics Service phục vụ thống kê.

Trong quá trình thực hiện, nhóm đã nghiên cứu và áp dụng nhiều công nghệ hiện đại như FastAPI, PostgreSQL, Docker Compose, OpenAPI Specification, RESTful API cùng cơ chế xác thực Bearer Token, Retry Mechanism và Idempotency nhằm đảm bảo hệ thống hoạt động ổn định, an toàn và dễ dàng mở rộng.

Báo cáo này trình bày toàn bộ quá trình nghiên cứu, phân tích, thiết kế, triển khai, kiểm thử và đánh giá Notification Service. Nội dung báo cáo được chia thành bốn chương, bao gồm tổng quan hệ thống, phân tích và thiết kế, triển khai và kiểm thử, kết quả thực nghiệm cùng đánh giá hệ thống. Thông qua báo cáo, nhóm mong muốn làm rõ quá trình xây dựng Notification Service cũng như những kết quả đạt được trong quá trình tích hợp với các nhóm phát triển khác trong hệ thống Smart Campus.
